%% ch09 --- Services: FastAPI and pydantic for people who know Minimal APIs
%%
%% Written September 2026 against the pinned versions. Every claim about the
%% framework in this chapter was checked against the installed package; the
%% three that contradicted the brief are recorded in CLAUDE.md under
%% *Resolved questions*.

\chapter{Services: FastAPI and pydantic for people who know Minimal APIs}
\label{ch:services}
\index{FastAPI}
\index{pydantic}

\noindent
A Minimal API endpoint is a function, its parameters are bound from the
request, and a container decides what the ones that are not part of the
request mean. FastAPI is the same idea with one substitution: the container
is the signature. There is no \code{IServiceCollection}, no registration, no
lifetime enum \dash{} there is a parameter with a default that says where its
value comes from, and pydantic is the model binder and the validator you used
to configure rather than to write.

The mapping is close enough that most of this chapter is a translation table.
It is worth your afternoon because of the three places the mapping is not
exact, and all three produce a service that starts, serves traffic and is
wrong: what a dependency's lifetime actually is, which end of the middleware
stack runs first, and what a worker buys you.

\mermaidfig{ch09-request-pipeline}{What a request meets, in order. Nothing here is configured; all of it is read off the signatures}{the request pipeline}

\section{The container is the signature}
\label{sec:ch09-depends}
\index{Depends@\texttt{Depends}}
\index{dependency injection}

Before reading the next listing, answer this from your .NET habit: a
dependency asked for twice in one request \dash{} is that one object or two,
and when is it disposed?

\pyregion{code/ch09/depends_scope.py}{scoped}{A dependency, and the two ways to ask for it}{lst:ch09-scoped}

\noindent
\code{Depends} takes a callable. FastAPI inspects the signature, calls it,
and passes the result in. A callable that \code{yield}s is a scope: what
follows the yield runs once the response has been sent, which is the
\code{using} block you would have written round a unit of work.

\pyregion{code/ch09/depends_scope.py}{route}{Two asks, one request}{lst:ch09-scoped-route}

\noindent
Run the file and the answer is not the one the habit predicts:

\transcript{ch09-depends}

\begin{trapbox}
\textbf{\code{Depends} is not a container, and its default is not a
singleton.} The default lifetime is per \emph{request}: ask for the same
callable in five parameters and it runs once, and the next request runs it
again. Transient is the thing you opt into, with
\code{Depends(fn, use\_cache=False)}. Singleton is not a lifetime at all
\dash{} it is something you build in the lifespan (\S\ref{sec:ch09-lifespan})
and reach through \code{request.state}. A registration that says
\code{AddSingleton} in your head and means \code{AddScoped} on the page is a
connection pool per request, and it will not fail until it is under load.
\end{trapbox}

\mermaidfig{ch09-depends-lifetime}{The three lifetimes, and which of them you have to ask for}{dependency lifetimes}

\begin{csbox}
\code{Annotated[Store, Depends(get\_store)]} is a type alias you can reuse,
which is the nearest thing here to registering the type once and injecting
it everywhere. Give it a name \dash{} \code{StoreDep} \dash{} and every
handler that needs a store takes one parameter of that alias. That is the
whole of the composition root.
\end{csbox}

\section{The binder is the model}
\label{sec:ch09-validation}
\index{pydantic!validation}
\index{HTTP status!422}

A parameter annotated with a pydantic model is the body; a parameter whose
name matches a path placeholder is from the path; anything else is a query
parameter. None of that is configured either.

\pyregion{code/ch09/service.py}{models}{The binder and the validator, in one declaration}{lst:ch09-models}

\noindent
What a .NET service does with a body it cannot bind is answer 400 and a
\code{ProblemDetails}. What this one does is answer 422, with pydantic's own
error list \dash{} one entry per failure, each naming its location, its
message and the input that caused it:

\transcript{ch09-validation}

\begin{trapbox}
\textbf{The status is 422 and the shape is pydantic's, not the framework's.}
There is no \code{ProblemDetails} and no
\code{InvalidModelStateResponseFactory}. Both halves matter: a caller written
against a .NET service will not recognise the status \emph{or} the body, and
a gateway that maps 4xx by hand will not have a rule for 422. If your callers
expect RFC 9457, one exception handler converts every validation failure in
the service at once, and the second half of the transcript above is what it
produces.
\end{trapbox}

\pyregion{code/ch09/validation.py}{problem}{One handler, and the whole service speaks problem details}{lst:ch09-problem}

\begin{exercise}{e09_02_problem_details}{Give the service the error shape your callers expect}
Write that handler yourself, against a test that checks the status, the
content type and the shape of the \code{errors} object. The interesting part
is not the mapping; it is that \code{exc.errors()} gives you a location
\emph{path} rather than a field name, and the first element of it is always
\code{body}.
\end{exercise}

\section{What goes out}
\label{sec:ch09-serialisation}
\index{pydantic!serialisation}
\index{OpenAPI}

\code{response\_model} is the serialisation contract, and it is a filter as
well as a shape: the handler may return a wider object and only the model's
fields reach the wire.

\pyregion{code/ch09/serialisation.py}{models}{Two models: what the service holds, and what the caller may see}{lst:ch09-out-models}

\pyregion{code/ch09/serialisation.py}{route}{The return annotation is the wide one; the response model is the narrow one}{lst:ch09-out-route}

\transcript{ch09-serialisation}

\noindent
Two things in that output are worth stopping at. The wire name is
\code{incidentId} while \code{model\_dump()} still gives you
\code{incident\_id}, because \code{serialization\_alias} renames the field
on the way out and nowhere else. And the generated schema carries the wire
name, so the OpenAPI document \dash{} which is produced from the signatures,
not from XML comments, and served at \code{/openapi.json} without a package
being added \dash{} agrees with what the service actually sends.

\begin{trapbox}
\textbf{\code{alias} is not \code{JsonPropertyName}: it renames the field on
the way IN as well.} A model whose field carries \code{alias="incidentId"}
can no longer be constructed as \code{Model(incident\_id=...)} \dash{} the
call raises a validation error naming a missing \code{incidentId}, and a
type checker says so before the test does. Use \code{serialization\_alias}
when you mean the wire format and \code{alias} only when the input really
does arrive under that name. When a model's input and output shapes differ,
FastAPI emits two schemas, suffixed \code{-Input} and \code{-Output}, where
Swashbuckle would have emitted one.
\end{trapbox}

\begin{exercise}{e09_03_response_model}{The wire name and the field name are different things}
Build a view model that renames two fields for the wire, keeps a secret out
of the response altogether, and is still constructible from Python with
ordinary keyword arguments. The test checks all three, and the fourth
assertion checks the generated schema, because a service whose spec and whose
output disagree is worse than one with no spec.
\end{exercise}

\section{Middleware, in the other order}
\label{sec:ch09-middleware}
\index{middleware!order}

Register two pieces of middleware, \code{audit} and then \code{timing}.
Before you read on, write down which one you expect to see the request first.

\pyregion{code/ch09/middleware_order.py}{order}{Two pieces of middleware, registered in this order}{lst:ch09-middleware}

\transcript{ch09-middleware}

\begin{trapbox}
\textbf{The last middleware registered is the outermost one.} In ASP.NET
Core the pipeline runs in registration order and the first \code{app.Use} is
the outer one; here it is the other way round, and the reason is one line of
Starlette: \code{add\_middleware} does
\code{self.user\_middleware.insert(0, ...)}, so every registration goes on
the front of the stack. Nothing warns you. The two orders are
indistinguishable until something depends on the sequence \dash{} an audit
log that must record what an authentication layer decided, say, or a timer
that is meant to include the work of everything inside it.
\end{trapbox}

\section{Configuration}
\label{sec:ch09-settings}
\index{pydantic-settings}
\index{configuration}

\pyregion{code/ch09/settings.py}{settings}{IOptions<T>, except that the binding is also the validation}{lst:ch09-settings}

\noindent
There is no provider chain, no \code{appsettings.json}, no
\code{appsettings.Production.json} layered over it and no
\code{IOptionsMonitor}. There is a model, its sources are the environment, a
\code{.env} file and a directory of secret files, and it is validated when it
is constructed \dash{} which, if you build it in the lifespan, means the
process refuses to start rather than failing at the first request that reads
a setting.

\begin{trapbox}
\textbf{\code{extra="forbid"} does not police the environment.} It rejects an
unknown key in the \code{.env} file and an unknown keyword passed to the
constructor, and it ignores an unknown \code{OPS\_} variable in the
environment entirely, because the environment source only ever looks up the
names the model declares. Measured at the pinned version, not assumed. The
prefix is real isolation in the other direction, though: a bare
\code{DATABASE\_URL} set by something else on the machine does not reach a
model declared with \code{env\_prefix="OPS\_"}.
\end{trapbox}

\begin{note}
A secret file's \emph{name} carries the prefix too, and the match is
case-insensitive: with \code{env\_prefix="OPS\_"}, a container secret
mounted as \code{api\_key} is not found and one mounted as
\code{ops\_api\_key} or \code{OPS\_API\_KEY} is. Declare the field as a
\code{SecretStr} and it stays out of every \code{repr} and every log line
until something calls \code{get\_secret\_value()}.
\end{note}

\begin{exercise}{e09_04_settings}{Configuration that fails at startup, not at first use}
Write the settings model for this service: a required database URL, a
bounded timeout, a worker count, and an API key read from a secrets
directory. Six tests, and two of them are about what the model must
\emph{refuse}.
\end{exercise}

\section{Lifespan, and work after the response}
\label{sec:ch09-lifespan}
\index{lifespan}
\index{background tasks}

\pyregion{code/ch09/service.py}{lifespan}{IHostedService's StartAsync and StopAsync, as one function}{lst:ch09-lifespan}

\noindent
Everything before the \code{yield} is startup and everything after it is
shutdown, which makes the pairing structural rather than a convention two
methods have to keep. What the lifespan yields becomes \code{request.state},
so the objects with process lifetime \dash{} a connection pool, a settings
object, an \code{httpx.AsyncClient} \dash{} are built once, in one place,
and are reachable from any handler.

For work that must happen after a response has been sent but does not deserve
a queue, a handler can take a \code{BackgroundTasks} parameter and add to it;
the tasks run after the response and after the middleware stack has unwound.
That is the \code{IHostedService} you would not have written: it is fine for
sending a notification and wrong for anything you would be sorry to lose,
because nothing here survives the process.

\section{Kestrel, uvicorn, and what a worker buys}
\label{sec:ch09-workers}
\index{uvicorn}
\index{measurement!E5}

uvicorn is the server and the application is an object it calls. The obvious
question is the one Kestrel answers with a thread pool: how do you use more
than one core? The answer is \code{--workers}, and before you read on, write
down what you expect four workers to do to the throughput of a service whose
handler spends its time awaiting something else.

Experiment E5 asks exactly that. One service, two handlers \dash{} one that
awaits \val{e05.delay.ms}~ms, standing in for an upstream call, and one that
burns \val{e05.cpu.ms}~ms of bytecode \dash{} driven over a real socket from
a separate process at a concurrency of \val{e05.conc.hi}, on a machine with
\val{e05.cores} cores. Every cell is measured against a calibration cell at
the same concurrency: the same driver against a handler that returns at once,
which reached \val{e05.noop.rps} requests a second, so nothing below is a
measurement of the client.

\begin{center}
\small
\begin{tabularx}{\linewidth}{@{}Xrrr@{}}
\toprule
 & \textbf{1 worker} & \textbf{4 workers} & \textbf{ratio} \\
\midrule
awaited upstream, requests/s & \val{e05.await.w1.rps} & \val{e05.await.w4.rps} & \val{e05.await.ratio4} \\
\quad p95 latency, ms & \val{e05.await.w1.p95} & \val{e05.await.w4.p95} & \\
CPU-bound, requests/s & \val{e05.cpu.w1.rps} & \val{e05.cpu.w4.rps} & \val{e05.cpu.ratio4} \\
\quad p95 latency, ms & \val{e05.cpu.w1.p95} & \val{e05.cpu.w4.p95} & \\
memory, MiB & \val{e05.pss.w1} & \val{e05.pss.w4} & \val{e05.pss.ratio4} \\
\bottomrule
\end{tabularx}
\end{center}

\noindent
The awaited row is the one to read first, and it is flat. One worker already
serves \val{e05.await.w1.rps} requests a second with a p50 of
\val{e05.await.w1.p50}~ms, which is the upstream delay and nothing else; at a
concurrency of one the same figure is \val{e05.await.w1.c1.p50}~ms. Four
workers move it by a factor of \val{e05.await.ratio4} and cost
\val{e05.pss.ratio4} times the memory. One loop holds forty awaited calls
without effort, because none of them is running.

\begin{trapbox}
\textbf{More workers is not more throughput; it is more processes.} A worker
is a process with its own loop, its own pools and its own memory, and it buys
exactly one thing: parallelism for work that holds the interpreter. For a
service that awaits, it buys nothing and costs \val{e05.pss.w4}~MiB where one
worker cost \val{e05.pss.w1}. Reach for it when a handler computes, and
reach for \code{asyncio.to\_thread} or a process pool first, because those do
not multiply the whole process.
\end{trapbox}

\noindent
The CPU row is the surprise: four workers on four cores moved it by a factor
of \val{e05.cpu.ratio4}, which is nothing at all. The reason is not the GIL
\dash{} these are separate processes \dash{} and the measurement says what it
is. Workers share one listening socket, and the worker that is in
\code{accept()} first takes the connection and keeps it: of
\val{e05.dist.conns} connections, only \val{e05.dist.reached} workers got any
at all and one of them took \val{e05.dist.busiest}. Every request on a
connection is answered by the same process for that connection's whole life.

\mermaidfig{ch09-worker-socket}{Why the CPU row did not scale: balancing happens once, at accept}{workers and the listening socket}

\begin{csbox}
Kestrel is one process with a thread pool, so a connection's requests can be
served by any thread and CPU-bound work spreads by itself. uvicorn's unit of
balancing is the \emph{connection}, not the request, so a few long-lived
keep-alive clients can pin most of your traffic onto one worker while the
others idle. That is an argument for a reverse proxy that opens many upstream
connections, and against assuming \code{--workers 4} has divided the work by
four.
\end{csbox}

\begin{versionbox}
\textbf{And \code{--workers} carries a transport defect at the pinned
version.} On the same machine, one keep-alive request answered in
\val{e05.nagle.w1}~ms with one worker and \val{e05.nagle.w4}~ms with four.
The cause is precise: \code{--workers} makes uvicorn bind the listening
socket itself, with \code{socket.socket(family=AF\_INET)}, whose
\code{proto} is 0 rather than \code{IPPROTO\_TCP}; accepted sockets inherit
that, and asyncio only sets \code{TCP\_NODELAY} on a socket whose
\code{proto} is \code{IPPROTO\_TCP}. So Nagle stays on, the response body is
a second write, and it waits out the client's delayed ACK. Forcing the client
to acknowledge at once takes the same exchange back to
\val{e05.nagle.w4.quick}~ms. An ordinary client does not do that, which is
why the same four-worker service measured \val{e05.plain.await.w4.rps}
requests a second at a p50 of \val{e05.plain.await.w4.p50}~ms against one
worker's \val{e05.plain.await.w1.rps}. Check it on your own build before you
blame your handler, and put a reverse proxy in front.
\end{versionbox}

\section{Testing, without a port}
\label{sec:ch09-testing}
\index{testing!TestClient}
\index{httpx!ASGITransport}

Neither of these starts a server, and that is the point: there is no port to
allocate, no readiness to wait for and nothing to leave running when a test
fails.

\pyregion{code/ch09/testing_client.py}{sync}{A WebApplicationFactory test, without the factory}{lst:ch09-testclient}

\pyregion{code/ch09/testing_client.py}{async}{The same test, on the loop the handler runs on}{lst:ch09-asgi}

\noindent
\code{TestClient} is synchronous and drives the application through a worker
thread, which is what you want for most tests. \code{ASGITransport} puts
\code{httpx.AsyncClient} straight on top of the application object, which is
what you want when the test is itself async \dash{} and it is the same client
class you would use against a real service, so a test and a caller share an
API.

\begin{exercise}{e09_05_async_client}{Exercise the application without starting a server}
Write a function that calls the application over \code{ASGITransport} and
returns the decoded body. One of its tests forbids \code{socket.socket}
outright, which is how you prove nothing went near the network.
\end{exercise}

\section{The skeleton, and the table}
\label{sec:ch09-skeleton}

Everything above is one file, in the order a request meets it: settings,
lifespan, models, a scoped dependency, two routes.

\pyfile{code/ch09/service.py}{The service skeleton: copy this, then delete what you do not need}{lst:ch09-service}

\begin{center}
\small
\begin{tabularx}{\linewidth}{@{}lX@{}}
\toprule
\textbf{ASP.NET Core} & \textbf{FastAPI} \\
\midrule
\code{app.MapGet("/x", handler)} & \code{@app.get("/x")} on the handler \\
\code{IServiceCollection} registration & a parameter with \code{Depends(fn)} \\
\code{AddScoped} & the default \\
\code{AddTransient} & \code{Depends(fn, use\_cache=False)} \\
\code{AddSingleton} & built in the lifespan, read from \code{request.state} \\
\code{IDisposable} on a scoped service & a dependency that \code{yield}s \\
model binding and \code{[ApiController]} & a pydantic model as a parameter \\
\code{ProblemDetails}, 400 & pydantic's error list, 422 \\
\code{JsonPropertyName} & \code{Field(serialization\_alias=...)} \\
\code{[JsonIgnore]} on a DTO & a narrower \code{response\_model} \\
Swashbuckle & \code{/openapi.json}, generated, no package \\
\code{app.Use(...)}, first is outermost & \code{@app.middleware("http")}, \textbf{last} is outermost \\
\code{IOptions<T>}, \code{appsettings.json} & \code{BaseSettings}, the environment and secret files \\
\code{IHostedService} & the \code{lifespan} context manager \\
Kestrel & uvicorn \\
threads in one process & \code{--workers}, which are processes \\
\code{WebApplicationFactory} & \code{TestClient} or \code{ASGITransport} \\
\bottomrule
\end{tabularx}
\end{center}

\begin{exercise}{e09_01_scoped_dependency}{A dependency with a lifetime}
Build the application the trap in \S\ref{sec:ch09-depends} describes: one
route that asks for a unit of work twice, one instance per request, a new one
next time, and closed when the request is over. Four tests, and the third is
the one that catches a dependency that returns instead of yielding.
\end{exercise}

\begin{note}
\textbf{What this chapter hands on.} Streaming a response \dash{} server-sent
events, a disconnecting client, a per-conversation lock \dash{} belongs to
the LangChain volume, which serves agents over exactly this stack and
measures the same service under load. The database session behind
\code{get\_store} is Chapter~\ref{ch:data}. What is here is the framework
mapping, and it is the same mapping whatever you serve with it.
\end{note}
